% Source: https://artofproblemsolving.com/community/c3h1994269_floor_function_problem
% Stable identity: aops-topic-1994269-q01
Let $n=\lfloor5x\rfloor$, an integer, so $3x^2=n$ and $n\le5x<n+1$. Since $3x^2\ge0$, $n\ge0$.

If $x<0$ then $5x<0$, so $n\le-1$, a contradiction. So $x\ge0$ and $x=\sqrt{\dfrac {n}{3}}$. For $n=0$ we get $x=0$, which works. For $n\ge1$ we need
\[
n\le5\sqrt{\dfrac {n}{3}}<n+1 .
\]
Squaring (all quantities are positive), the left inequality is $n^2\le\dfrac{25n}{3}$, i.e.\ $n\le8$. The right inequality is $\dfrac{25n}{3}<(n+1)^2$, i.e.\ $3n^2-19n+3>0$; the roots of $3n^2-19n+3$ are $\dfrac{19\pm\sqrt{325}}{6}\approx0.16$ and $6.17$, so among the integers $1\le n\le8$ this holds only for $n=7$ and $n=8$.

Check: $x=\sqrt{\dfrac{7}{3}}\approx1.528$ gives $5x\approx7.64$, floor $7=3x^2$; and $x=\sqrt{\dfrac{8}{3}}\approx1.633$ gives $5x\approx8.16$, floor $8=3x^2$. So the solutions are
\[
\boxed{x=0,\quad x=\dfrac{\sqrt{21}}{3},\quad x=\dfrac{2\sqrt6}{3}}.
\]
(The value $x=\dfrac{5}{3}$ is \emph{not} a solution: $3x^2=\dfrac{25}{3}$ is not an integer.)
